\documentclass[../../main.tex]{subfiles}
\begin{document}
\section{直积与直和}\label{section:appB-direct-products-and-sums}

如果 $V_1,\dots,V_k$ 是实向量空间，它们的\emph{直积}是底集合为 Cartesian 积 $V_1\times\cdots\times V_k$ 的向量空间，加法与数乘按分量定义：
\begin{align*}
(v_1,\dots,v_k)+(v'_1,\dots,v'_k)=(v_1+v'_1,\dots,v_k+v'_k),
\end{align*}
\begin{align*}
c(v_1,\dots,v_k)=(cv_1,\dots,cv_k).
\end{align*}
基本例子是 Euclid 空间 $\mathbb R^n=\mathbb R\times\cdots\times\mathbb R$。

对某些应用（主要是我们在第 17、18 章对 de Rham 上同调的处理），把这推广到无限多个向量空间很重要。为作此讨论，我们转向交换环 $R$ 上模的一般背景。$R$-模之间的线性映射恰如向量空间那样定义：如果 $V,W$ 是 $R$-模，映射 $F:V\to W$ 如果对所有 $a,b\in R$ 与 $v,w\in V$ 都有 $F(av+bw)=aF(v)+bF(w)$，就称它是 $R$-线性的。如果 $V$ 是 $R$-模，子集 $S\subset V$ 如果在加法与数乘下封闭，就称为 $V$ 的\emph{子模}，于是它本身也是 $R$-模。本节余下部分假设 $R$ 是固定的交换环。在我们所有应用中，$R$ 要么是实数域 $\mathbb R$（此时模是实向量空间、线性映射是通常的线性映射），要么是整数环 $\mathbb Z$（此时模是 Abel 群、线性映射是群同态）。

如果 $(V_\alpha)_{\alpha\in A}$ 是任意下标化集合族，它们的 Cartesian 积记作 $\prod_{\alpha\in A}V_\alpha$，定义为函数 $v:A\to\bigcup_{\alpha\in A}V_\alpha$ 所成集合，满足对每个 $\alpha$ 有 $v(\alpha)\in V_\alpha$。由选择公理，非空下标的非空集合族的 Cartesian 积非空。如果 $v$ 是 Cartesian 积的元素，我们通常把 $v$ 在 $\alpha\in A$ 处的值记作 $v_\alpha$ 而非 $v(\alpha)$；元素 $v$ 本身通常记作 $(v_\alpha)_{\alpha\in A}$，或当下标集已知时简记 $(v_\alpha)$。这可看作集合 $V_\alpha$ 元素的一个下标族，或一个“$A$-元组”。对每个 $\beta\in A$，有一个典范投影映射 $\pi_\beta:\prod_{\alpha\in A}V_\alpha\to V_\beta$，定义为
\begin{align*}
\pi_\beta((v_\alpha)_{\alpha\in A})=v_\beta.
\end{align*}

现设 $(V_\alpha)_{\alpha\in A}$ 是一个 $R$-模下标族。该族的\emph{直积}是集合 $\prod_{\alpha\in A}V_\alpha$，通过如下定义加法与数乘而成为 $R$-模：
\begin{align*}
(v_\alpha)+(v'_\alpha)=(v_\alpha+v'_\alpha),\qquad c(v_\alpha)=(cv_\alpha).
\end{align*}
这个模的零元素是对每个 $\alpha$ 都有 $v_\alpha=0$ 的 $A$-元组。容易验证每个投影映射 $\pi_\beta$ 都是 $R$-线性的。

\begin{proposition}[直积的特征性质]\label{proposition:appB-direct-product-universal}
设 $(V_\alpha)_{\alpha\in A}$ 是 $R$-模下标族。给定一个 $R$-模 $W$ 和一族 $R$-线性映射 $G_\alpha:W\to V_\alpha$，存在唯一的 $R$-线性映射 $G:W\to\prod_{\alpha\in A}V_\alpha$，使对每个 $\alpha\in A$ 有 $\pi_\alpha\circ G=G_\alpha$。
\end{proposition}

\begin{exercise}\label{exercise:appB-B-54}
证明前一命题。
\end{exercise}

与直积互补的是直和的概念。给定如上下标族 $(V_\alpha)_{\alpha\in A}$，我们把该族的\emph{直和}定义为它们直积的子模，由满足 $v_\alpha=0$（除至多有限个 $\alpha$ 外）的 $A$-元组 $(v_\alpha)_{\alpha\in A}$ 组成。直和记作 $\bigoplus_{\alpha\in A}V_\alpha$；对有限族则记作 $V_1\oplus\cdots\oplus V_k$。对有限模族，直积与直和相同。

对每个 $\beta\in A$，有一个典范 $R$-线性单射 $\iota_\beta:V_\beta\to\bigoplus_{\alpha\in A}V_\alpha$，它把 $\iota_\beta(v)$ 定义为 $A$-元组 $(v_\alpha)_{\alpha\in A}$，其中 $v_\beta=v$ 而对 $\alpha\neq\beta$ 有 $v_\alpha=0$。在有限直和的情形，这恰表示 $\iota_\beta(v)=(0,\dots,0,v,0,\dots,0)$，其中 $v$ 在第 $\beta$ 个位置。

\begin{proposition}[直和的特征性质]\label{proposition:appB-direct-sum-universal}
设 $(V_\alpha)_{\alpha\in A}$ 是 $R$-模下标族。给定一个 $R$-模 $W$ 和一族 $R$-线性映射 $G_\alpha:V_\alpha\to W$，存在唯一的 $R$-线性映射 $G:\bigoplus_{\alpha\in A}V_\alpha\to W$，使对每个 $\alpha\in A$ 有 $G\circ\iota_\alpha=G_\alpha$。
\end{proposition}

\begin{exercise}\label{exercise:appB-B-56}
证明前一命题。
\end{exercise}

如果 $W$ 是 $R$-模，$(V_\alpha)_{\alpha\in A}$ 是 $W$ 的一族子空间，则把特征性质用于包含映射 $\iota_\alpha:V_\alpha\hookrightarrow W$，就保证存在一个典范 $R$-线性映射 $\bigoplus_\alpha V_\alpha\to W$，它在每个 $V_\alpha$ 上限制为包含。这个映射是同构，当且仅当诸 $V_\alpha$ 选取得使 $W$ 的每个元素都有唯一方式表为有限线性组合 $\sum_\alpha c_\alpha v_\alpha$（其中对每个 $\alpha$ 有 $v_\alpha\in V_\alpha$）。在此情形中，我们可自然地把 $W$ 等同于 $\bigoplus_\alpha V_\alpha$，并称 $W$ 是子模族 $\{V_\alpha\}$ 的\emph{内直和}，把我们此前对向量空间两个互补子空间引入的术语加以推广。抽象模族的直和有时称为\emph{外直和}，以区别于内直和。

如果 $V,W$ 是 $R$-模，从 $V$ 到 $W$ 的所有 $R$-线性映射所成集合 $\operatorname{Hom}_R(V;W)$ 在逐点加法与数乘下是 $R$-模：
\begin{align*}
(F+G)(v)=F(v)+G(v),\qquad (aF)(v)=aF(v).
\end{align*}
（如果 $V,W$ 是实向量空间，则 $\operatorname{Hom}_R(V;W)$ 正是 $\mathbb R$-线性映射空间 $\mathcal L(V;W)$；如果它们是 Abel 群，则 $\mathbb Z$-线性映射就是群同态，我们通常写 $\operatorname{Hom}(V;W)$ 而非 $\operatorname{Hom}_{\mathbb Z}(V;W)$。）我们最后一个命题用于第 18 章 de Rham 定理的证明。

\begin{proposition}\label{proposition:appB-hom-direct-sum}
设 $(V_\alpha)_{\alpha\in A}$ 是 $R$-模下标族。对每个 $R$-模 $W$，存在典范同构
\begin{align*}
\operatorname{Hom}_R\left(\bigoplus_{\alpha\in A}V_\alpha;W\right)\cong\prod_{\alpha\in A}\operatorname{Hom}_R(V_\alpha;W).
\end{align*}
\end{proposition}

\begin{proof}
定义映射 $\Phi:\operatorname{Hom}_R\left(\bigoplus_{\alpha\in A}V_\alpha;W\right)\to\prod_{\alpha\in A}\operatorname{Hom}_R(V_\alpha;W)$，令 $\Phi(F)=(F_\alpha)_{\alpha\in A}$，其中 $F_\alpha=F\circ\iota_\alpha$。

为证 $\Phi$ 满射，设 $(F_\alpha)_{\alpha\in A}$ 是 $\prod_{\alpha\in A}\operatorname{Hom}_R(V_\alpha;W)$ 的任意元素。这意味着对每个 $\alpha$，$F_\alpha$ 是从 $V_\alpha$ 到 $W$ 的 $R$-线性映射。直和的特征性质于是保证存在一个 $R$-线性映射 $F:\bigoplus_{\alpha\in A}V_\alpha\to W$，满足对每个 $\alpha$ 有 $F\circ\iota_\alpha=F_\alpha$，这等价于 $\Phi(F)=(F_\alpha)_{\alpha\in A}$。

为证 $\Phi$ 单射，设 $\Phi(F)=(F_\alpha)_{\alpha\in A}=0$。由直积零元素的定义，这意味着对每个 $\alpha$，$F_\alpha=F\circ\iota_\alpha$ 都是零同态。由命题 B.55 的唯一性断言，这蕴含 $F$ 本身是零同态。
\end{proof}
\end{document}
